\chapter*{English Abstract}
\addcontentsline{toc}{chapter}{English Abstract}

This Master's thesis investigates the design and prototype implementation of a telemedicine platform (Shifa) for the healthcare context of Uzbekistan. The work is grounded in qualitative physician findings that indicate recurring delivery problems, including informal and fragmented communication, manual appointment coordination, discontinuous documentation practices, exposure to data-loss risks, and legal uncertainty in remote-care settings.

Based on these findings, the thesis derives functional and technical requirements and translates them into an implementable system architecture. The resulting prototype consists of two client applications (patient and doctor/admin), a central backend with role-aware access control, and PostgreSQL-based persistence. External services are integrated for OTP-oriented authentication, notifications, and video consultation workflows.

The contribution of the thesis is the traceable link between empirical needs and concrete implementation decisions. At the same time, the document explicitly reports prototype boundaries, especially regarding production-grade hardening, operational maturity, and fully formalized API contracts.
